\documentclass[10pt,letterpaper]{article}

% ---------- Page setup ----------
% Prism-compatible version: compile with pdfLaTeX.
% Uses Helvetica/Arial-like sans-serif throughout because many cloud LaTeX engines
% do not expose the Microsoft Arial system font.
\usepackage[margin=0.52in]{geometry}
\usepackage[T1]{fontenc}
\usepackage[utf8]{inputenc}
\usepackage[scaled=0.95]{helvet}
\renewcommand{\familydefault}{\sfdefault}
\usepackage{microtype}
\usepackage{tabularx}
\usepackage{enumitem}
\usepackage{titlesec}
\usepackage[hidelinks]{hyperref}
\usepackage{xcolor}
\usepackage{ragged2e}

\pagestyle{empty}
\setlength{\parindent}{0pt}
\setlength{\parskip}{0.6pt}
\setlist[itemize]{leftmargin=1.05em, itemsep=0.7pt, topsep=0.7pt, parsep=0pt}
\urlstyle{same}
\sloppy

% ---------- Section style ----------
\definecolor{sectiongray}{gray}{0.18}
\titleformat{\section}
  {\normalsize\bfseries\color{sectiongray}}
  {}{0pt}{}
  [\vspace{1.8pt}{\color{sectiongray}\titlerule[0.5pt]}]
\titlespacing*{\section}{0pt}{7pt}{5pt}

% ---------- Custom commands ----------
\newcommand{\cvsubheading}[4]{%
  \begin{tabularx}{\textwidth}{@{}X r@{}}
    \textbf{#1} & #3 \\
    \textit{#2} & \textit{#4} \\
  \end{tabularx}\vspace{-3pt}
}
\newcommand{\publication}[1]{\item \justifying #1}

\begin{document}
\fontsize{9.25}{10.25}\selectfont

% ---------- Header ----------
\begin{center}
    {\Large \textbf{Qianyi Liu}} \\
    \vspace{1pt}
    Ph.D. Student, Department of Civil Engineering, Stony Brook University \\
    \vspace{1pt}
    \href{mailto:qianyi.liu@stonybrook.edu}{qianyi.liu@stonybrook.edu} \; | \; Stony Brook, NY, USA \; | \; ASCE Student Member \; | \; AGU Student Member
\end{center}
\vspace{-6pt}

\section*{EDUCATION}
\cvsubheading{Stony Brook University}{Ph.D. in Civil Engineering; GPA: 4.00/4.00}{08/2025--Present}{Stony Brook, NY, USA}
\begin{itemize}
    \item Research focus: geotechnical/geohazard engineering, landslide susceptibility mapping, GIS, spatial machine learning, causal and explainable AI. Advisor: Prof. Te Pei.
\end{itemize}

\cvsubheading{University of Macau}{M.S. in Coastal Environment and Safety; GPA: 3.25/4.00}{08/2023--06/2025}{Macau, China}
\begin{itemize}
    \item Relevant modules: Hydrodynamics and Floods (A); Advanced Computational Methods: Principles and Applications (A-).
\end{itemize}

\cvsubheading{Capital University of Economics and Business}{B.Eng. in Safety Engineering; GPA: 3.81/5.00 (88.2/100)}{09/2019--06/2023}{Beijing, China}
\begin{itemize}
    \item Relevant modules: Electrical Safety Engineering (100), Systems Engineering (100), Engineering Drawing II (99), Python Data Analysis (93).
\end{itemize}

\section*{RESEARCH INTERESTS}
Geotechnical engineering; landslide susceptibility mapping; geohazard assessment; trustworthy AI; causal and explainable machine learning; spatially aware machine learning; GIS and remote sensing; hydrodynamic and compound-flood modeling; spatial cross-validation; SHAP analysis; infrastructure resilience.

\section*{RESEARCH EXPERIENCE}
\cvsubheading{Landslide Susceptibility Mapping and Geohazard Assessment}{Ph.D. Researcher, Stony Brook University}{08/2025--Present}{Stony Brook, NY, USA}
\begin{itemize}
    \item Developing geospatial machine-learning workflows for landslide susceptibility and slope-related geohazard assessment, emphasizing physically meaningful, spatially robust, and interpretable outputs.
    \item Integrating terrain, hydrological, soil-related, vegetation, and rainfall-related factors to analyze spatially varying slope-instability mechanisms.
    \item Applying spatial cross-validation, explainable AI, causal-graph analysis, and GIS-based mapping to improve model reliability and engineering interpretability.
\end{itemize}

\cvsubheading{Flood Susceptibility Mapping with Explainable Ensemble Machine Learning}{Research Assistant, University of Macau}{08/2023--08/2024}{Macau, China}
\begin{itemize}
    \item Conducted research on flood susceptibility mapping using GIS, multi-source environmental data, hydrodynamic modeling, and explainable ensemble machine-learning methods.
    \item Developed flood susceptibility maps for the Greater Bay Area / Pearl River Delta case study and analyzed feature contributions using SHAP-based interpretation.
\end{itemize}

\cvsubheading{Smart Emergency Response Development Strategy (2035)}{National Key Project; Team Member; Lead: Academician Weicheng Fan}{09/2020--10/2020}{China}
\begin{itemize}
    \item Synthesized national macro policies related to smart emergency management and contributed to drafting the policy environment and national demand section for smart emergency development.
\end{itemize}

\section*{PUBLICATIONS AND MANUSCRIPTS}
\textbf{Peer-Reviewed Journal Articles}
\begin{itemize}
    \publication{\textbf{Liu, Q.}, Gao, L., Du, H., Pei, T., Kuok, S.-C., \& Chen, S. (2026). Integrating explainable ensemble machine learning with hydrodynamic modeling to assess tropical cyclone-induced flood susceptibility. \textit{Geoscience Frontiers}, 17(5), Article 102355. \href{https://doi.org/10.1016/j.gsf.2026.102355}{https://doi.org/10.1016/j.gsf.2026.102355}.}
    \publication{Ma, S., \textbf{Liu, Q.}, \& Zhang, Y. (2021). A prediction method of fire frequency: Based on the optimization of the SARIMA model. \textit{PLOS ONE}, 16(8), Article e0255857. \href{https://doi.org/10.1371/journal.pone.0255857}{https://doi.org/10.1371/journal.pone.0255857}.}
\end{itemize}
\textbf{Manuscripts Under Review}
\begin{itemize}
    \publication{\textbf{Liu, Q.}, Liu, X., Cui, H., Tian, Y., Devineni, N., \& Pei, T. (2026). \textit{Spatially Robust Explainable Machine Learning for Landslide Susceptibility Mapping: Interpreting Spatially Heterogeneous Controls}. Submitted to \textit{Bulletin of Engineering Geology and the Environment}. Status: Under Review.}
    \publication{Du, H., Fei, K., \textbf{Liu, Q.}, \& Zhang, L. (2026). \textit{Tropical Cyclone-Induced Compound Flooding Reveals Critical Flood Zone Transitions in the Pearl River Delta}. Submitted to \textit{Ocean Engineering}. Status: Under Review.}
\end{itemize}

\section*{CONFERENCE ABSTRACTS AND PRESENTATIONS}
\begin{itemize}
    \publication{\textbf{Liu, Q.}, Liu, X., Pei, T., Devineni, N., \& Tian, Y. (2025). Understanding Geo-Environmental Factors for Landslide Occurrences Using Geographically Explainable AI: A Case Study from the Northeastern United States. \textit{AGU Fall Meeting Abstracts}, 2025, NH33F-0502.}
    \publication{\textbf{Liu, Q.}, Liu, X., \& Pei, T. (2026). Uncovering Spatial Heterogeneity in Landslide Controls through a Causality-Aware and Explainable Machine Learning Framework. 2026 NSF NHERI GSC Mini-Conference.}
    \publication{\textbf{Liu, Q.}, Liu, X., \& Pei, T. (2026). Trustworthy AI for Geohazard Mitigation: A Causal and Explainable Framework for Landslide Susceptibility Mapping. SUNY AI Symposium.}
    \publication{\textbf{Liu, Q.}, Liu, X., Pei, T., Devineni, N., \& Tian, Y. (2027). Physics-Guided and Spatially Aware Machine Learning for Trustworthy Landslide Susceptibility Mapping in Heterogeneous Terrain. ASCE 2027 Abstract.}
    \publication{\textbf{Liu, Q.}, Liu, X., Pei, T., Cui, H., Devineni, N., \& Tian, Y. (2027). Causality-Constrained Explainable Machine Learning for Identifying Spatially Varying Controls in Landslide Susceptibility Mapping. ASCE 2027 Abstract.}
\end{itemize}

\newpage
\section*{TEACHING EXPERIENCE}
\cvsubheading{Department of Civil Engineering, Stony Brook University}{Teaching Assistant}{08/2025--Present}{Stony Brook, NY, USA}
\begin{itemize}
    \item Support course instruction, grading, laboratory/problem-solving sessions, and student guidance for undergraduate civil engineering courses.
    \item Help students connect experimental or computational results with civil engineering concepts and reporting standards.
\end{itemize}

\section*{ACADEMIC SERVICE}
\cvsubheading{2026 NSF NHERI GSC Mini-Conference}{Reviewer, NSF NHERI Graduate Student Council}{04/2026}{}
\begin{itemize}
    \item Reviewed student conference submissions and contributed to the academic review process for the 2026 NHERI GSC Mini-Conference.
\end{itemize}

\section*{ADDITIONAL PROFESSIONAL EXPERIENCE}
\cvsubheading{China Aerospace Construction Group Co., Ltd.}{Intern}{06/2022--07/2022; 09/2022--10/2022}{Beijing, China}
\begin{itemize}
    \item Compiled safety inspection records for the Pinggu Men Station natural gas project, organized ledgers, and developed educational safety materials and emergency response plans.
    \item Supported implementation of corporate safety protocols and enhanced execution methods for safety management under pandemic-related constraints.
\end{itemize}

\cvsubheading{Beijing GuoXinAn Technology Co., Ltd.}{Intern}{07/2021--08/2021}{Beijing, China}
\begin{itemize}
    \item Assisted SMEs with risk-source declaration in Changping District through direct offline engagement and online reporting support.
    \item Supported community-level enterprise safety and regional risk management practices.
\end{itemize}

\section*{AWARDS AND HONORS}
\begin{itemize}
    \item First Place, Individual Poster Presentation, 2026 NSF NHERI GSC Mini-Conference, 2026.
    \item Second Prize, National College Students' Mathematics Modelling Competition, Team Leader, 09/2021.
    \item Second Prize, National College Students' Mathematics Modelling Competition, 09/2020.
    \item Honorable Mention, Mathematical Contest in Modeling (MCM/ICM), 09/2021.
    \item Second Prize, Beijing Energy Conservation and Emission Reduction Competition, Team Leader, 09/2021.
    \item Key Project Leader, Summer Social Practice by the College, 09/2021.
\end{itemize}

\section*{PROFESSIONAL MEMBERSHIPS}
\begin{itemize}
    \item American Society of Civil Engineers (ASCE), Student Member.
    \item American Geophysical Union (AGU), Student Member.
\end{itemize}

\section*{TECHNICAL SKILLS}
\begin{itemize}
    \item \textbf{Geospatial analysis:} ArcGIS Pro, ArcMap, GIS-based spatial analysis, raster data processing, susceptibility mapping.
    \item \textbf{Programming and data analysis:} Python, Jupyter Notebook, MATLAB, statistical analysis, Microsoft Office Suite.
    \item \textbf{Machine learning and interpretation:} Ensemble machine learning, trustworthy AI, causal and explainable AI, SHAP analysis, model evaluation and visualization.
    \item \textbf{Specialized software:} SPSS, Stata, Citespace, Origin. \textbf{Languages:} Mandarin Chinese (native), English (fluent). \textbf{Certificates:} CET-4, CET-6.
\end{itemize}

\section*{SELECTED INTERESTS}
Swimming; model aviation; Pokemon Trading Card Game.

\end{document}
